% !TEX root = ../main.tex

\chapter{Métricas e Conexões Afins}
\label{chap:metricas_e_conexoes_afins}

O objetivo desse capítulo é o de explicitar a noção de métrica riemmaniana, a qual generaliza a primeira forma fundamental estudada em Geometria Diferencial. Aqui, deixamos de estudar topologia diferencial e passamos, de fato, a estudar Geometria Riemmanaiana, quando analisamos variedades riemmanas, as quais são variedades diferenciáveis munidas de uma métrica.

Uma métrica riemanniana em uma variedade diferenciável é uma regra para definir um produto interno definido positivo em cada espaço tangente, que varia diferenciavelmente em um sentido a ser formalizado. A grosso modo, uma métrica riemaniana permite calcular distâncias e ângulos em uma variedade diferenciável, assim como uma métrica em um espaço euclidiano.

Mostramos que toda variedade diferenciável bem comportada pode ser munida de uma métrica riemaniana, tornando-se uma variedade riemaniana. Esse resultado é interessante pois extende a aplicabilidade desse capítulo e dos posteriores a uma grande classe de variedades diferenciáveis. Posteriormente, quando apresentamos o teorema de Hopf-Rinow, vemos que tais ferramentas da Geometria Riemanniana nos ajudam a deduzir certas propriedades das próprias variedades diferenciavéis.

Também introduzimos o conceito de conexão, o qual serve como uma forma de identificar vetores em diferentes espaços tangentes, permitindo extender o conceito de diferenciabilidade ao longo de campos vetoriais. Mostramos o teorema de Levi-Civita, que afirma que toda variedade riemanniana possuí uma única conexão compatível com a topologia e com a métrica da variedade, em um sentido a ser tornado preciso.
\section{Métricas Riemannianas}
\label{loc:capítulo_3_:_métricas_e_conexões.métricas_riemannianas}
Iniciamos por formalizar a noção de métrica:
\begin{definition}
    \label{loc:2.1_:_def_1_:_métrica_riemanniana.definição}
    Seja $M$ uma variedade diferenciável. Uma métrica Riemanniana é uma correspondência que associa a cada $p \in M$, um produto interno $\left<\cdot, \cdot\right>_{p}$, isto é, uma aplicação bilinear $\left<\cdot,\cdot\right>_{p}:T_{p}M\times T_{p}M\to \mathbb{R}$ que varia diferencialmente com respeito a $p$ e que satisfaz:
    \begin{enumerate}
        \item Simétrica: $\left<u, v\right>_{p}= \left<v,u\right>_{p}$;
        \item Definida positiva: $\left<v, v\right>_{p} \geq 0$ e $\left<v,v\right>_{p}=0$ se, e somente se, $v=0$.
    \end{enumerate}

    Variar diferencialmente com respeito a $p$ significa que, dada uma parametrização $\mathbf{x}:U\to M$, em que $\mathbf{x}(x_{1}, \dots, x_{n})=q$ e $\frac{\partial}{\partial x_{i}}(q)=d\mathbf{x}_{q}(0, .., 0, 1, 0, \dots, 0)$, então $g_{ij}(x_{1}, \dots, x_{n})=\left<\frac{\partial}{\partial x_{i}}(q),\frac{\partial}{\partial x_{j}}(q)\right>_{q}$ é diferenciável em $U$.
\end{definition}
Dizemos que uma variedade diferenciável munida de uma métrica riemanniana é uma variedade riemanniana.

A transformação que preserva propriedades derivadas da métrica é denominada isometria, e é definida da seguinte maneira:
\begin{definition}
    \label{loc:definicao_isometria}
    $ $
    \begin{itemize}
        \item Sejam $M,N$ variedades diferenciáveis. Um difeomorfismo $f:M\to N$ é chamado \textbf{isometria} se
              \begin{equation*}
                  \left< u, v \right>_{p}=\left< df_{q}(u), df(q)(v) \right>_{p}, \text{para todo } u,v \in T_{p}M. \tag{1}
              \end{equation*}
        \item Dizemos que $f:M\to N$ é uma \textbf{isometria local} se, para todo $p \in M$, existe uma vizinhança $U$ de $p$ em $M$, tal que (1) é satisfeito para todo $q \in U$ e $u,v \in T_{q}M$.
        \item Dizemos que $M$ é \textbf{localmente isométrica} a $N$ se, para todo $p \in M$, existe uma vizinhança $U$ de $p$ em $M$ e uma isometria $f:U\to f(U) \subset N$, isto é, existe um difeomorfismo $f:U\to f(U)$, que satisfaz (1).
    \end{itemize}
\end{definition}
Note que $M$ localmente isométrica a $N$ não implica que $N$ é localmente isométrica a $M$. De fato: \textbf{PLACEHOLDER: CITAR EXEMPLO. É UM EXERCÍCIO DO LIVRO, AINDA NÃO CONSEGUI RESOLVER!}
\begin{examples}
    \label{loc:exemplos_de_métricas_e_isometrias.exemplos_simples}
    $ $
    \begin{enumerate}
        \item $M=\mathbb{R}^{n}$, com $\frac{\partial}{\partial x_{i}}$ identificado com $e_{i}$. Definimos uma métrica em  $M$, pondo

              \begin{equation*}
                  g_{ij}=\left< e_{i}, e_{j} \right>= \delta_{ij}= \begin{cases} 1, & \text{se } i=j,\\ 0, & \text{se } i\neq j. \end{cases}
              \end{equation*}
              A bilinearidade e simetria seguem trivialmente e o fato da métrica ser diferenciável segue do fato de ser constante, fixados $i,j$. Essa métrica é denominada de métrica canônica ou euclidiana.
        \item Variedades Imersas:
              Seja $f:M^{n}\to N^{n+k}$ uma imersão, isto é, $f$ é diferenciável e $df_{p}$ é injetora para todo $p \in M$.
              Se $N$ é uma variedade riemaniana, podemos induzir uma métrica em $M$, pondo $\left< u, v \right>_{p}=\left< df_{p}(u), df_{p}(v) \right>_{f(p)}$. De fato, isso define uma métrica:
              \begin{itemize}
                  \item A bilinearidade segue do fato de $df_{p}$ ser linear;
                  \item A simetria segue da métrica de $N$;
                  \item Para mostrar que é positiva definida, tome $v \in T_{p}M$. Temos então que $\left< v, v \right>_{p}=\left< df_{p}(v), df_{p}(v) \right>f_{p}\geq0$. No caso se ser igual a $0$, temos que $df_{p}(v)=0$ e, como $df_{p}$ é injetiva, $v=0$.
                        Note que precisamos da injetividade de $df_{p}$ apenas para mostrar o fato de que a métrica é definida positiva.
                  \item Podemos juntar os exemplos 1. e 2; para construir uma métrica na esfera $S^{n}$. Com efeito, a função inclusão $i:S^{n}\to \mathbb{R}^{n+1}$ é um mergulho, em particular, é uma imersão. Logo, munimos $\mathbb{R}^{n+1}$ da métrica canônica e $S^{n}$ da métrica induzida por tal.
                  \item Considere o semiplano superior aberto $\mathbb{R}^{2}_{+}=\{ (x,y) \in \mathbb{R}^{2}; y >0 \}$ munido da estrutura diferencial $\{ (\mathbb{R}^{2}_{+},Id) \}$ com a métrica $\left<u,v\right>_{(x,y)}= \frac{1}{y}\left<u,v\right>_{\mathbb{R}^{2}}$, onde $\left<u,v\right>_{\mathbb{R}^{2}}$ denota a métrica euclidiana em $\mathbb{R}^{2}$. Essa variedade riemannaina é chamada de modelo de semiplano de poincaré e ela serve como modelo da geometria hiperbólica de Lobachevsky.  \textbf{PLACEHOLDER: REFERÊNCIA}.
                  \item Análogamente ao semiplano de poincáre, temos o disco de poincáre, o qual é construido tomando $D=\{ (x,y)\in \mathbb{R}^{2}|x^{2}+y^{2}< 1 \}$, munindo-o com a estrutura diferencial identidade e com a métrica $\left<u,v\right>_{(x,y)}=\frac{4\left<u,v\right>_{\mathbb{R}^{2}}}{(1-(x^{2}+y^{2}))^{2}}$, temos outro modelo da geometria hiperbólica. A idealização dessa variedade é encontrada no livro Ciência e Hipótese \cite{poincare1905}, onde a métrica riemanniana é substituida pela dilatação térmica. Poincaré explica intuitivamente da seguinte maneira:
                        \begin{citacao}
                            Observemos, em primeiro lugar, que embora do ponto de vista da nossa geometria comum este mundo seja finito, para os seus habitantes ele parecerá infinito. À medida que se aproximam da superfície da esfera, eles tornam-se mais frios e, ao mesmo tempo, cada vez menores. Os passos que dão são, portanto, também cada vez menores, de modo que nunca poderão alcançar a fronteira da esfera (p.65)$\footnote{Traduzido pelos autores.}$.
                        \end{citacao}
              \end{itemize}
    \end{enumerate}
\end{examples}
\begin{example}
    \label{loc:exemplos_de_métricas_e_isometrias.métrica_produto}
    Sejam $M_{1}$ e $M_{2}$ variedades riemannianas. Definimos uma métrica na variedade produto $M_{1}\times M_{2}$ pondo
    \begin{equation*}
        \left< u, v \right>_{(p,q)}=\left< d\pi_{1}(u), d\pi_{1}(v) \right>_{p}+\left< d\pi_{2}(u),d\pi_{2}(v)  \right>_{q}
    \end{equation*}
    onde $\pi_{1}:M_{1}\times M_{2} \to M_{1}$ e $\pi_{2}:M_{1}\times M_{2}\to M_{2}$ são as projeções.
\end{example}
O segundo item do \Cref{loc:exemplos_de_métricas_e_isometrias.exemplos_simples} será utilizado posteriormente, quando estudarmos o mapa exponencial.

É interessante observar que, na Geometria Diferencial, a primeira forma fundamental é um exemplo de métrica riemanniana, a qual pode ser obtida definindo em $\mathbb{R}^{3}$ a métrica euclidiana e induzindo tal métrica numa superfície, a qual está imersa em $\mathbb{R}^{3}$.

\textbf{PLACEHOLDER: AQUI EXISTE A POSSIBILIDADE DE COLOCAR O EXEMPLO GIGANTESCO DE GRUPOS DE LIE. CREIO NÃO SER NECESSÁRIO POR QUE NÃO CONTRIBUI NADA PARA O OBJETIVO DO TCC}

Agora, generalizamos a noção de curva diferenciável para a de curva parametrizada.
\begin{definition}
    \label{loc:2.8_:_def_3_:_curva_parametrizada.definição}
    $ $
    \begin{itemize}
        \item \textbf{Curva parametrizada:} Uma aplicação diferenciável $c:I\to M$, onde $I \subset \mathbb{R}$ é um intervalo aberto é chamada de curva parametrizada.
        \item \textbf{Segmento:} Se $a,b \in I$ e $a <b$, dizemos que $c|_{[a,b]}$ é um segmento da curva $c$. Também nos referiremos a $c([a,b])\subset M$ como segmento.
    \end{itemize}
\end{definition}
De agora em diante, utilizaremos o termo curva diferenciável tanto para se referir à aplicações $c:(-\epsilon, \epsilon)\to M$ diferenciáveis quanto à curvas parametrizadas.

Agora, dada uma curva diferenciável, buscamos análisar sua velocidade e, posteriormente, sua aceleração. Na Geometria Diferencial, isso é feito considerando a curva como uma curva em $\mathbb{R}^{3}$ e tomando sua derivada usual. Na Geometria Riemannana, como não dispomos de um espaço ambiente euclidiano, definimos a velocidade observando os espaços tangentes. Primeiramente, de forma mais geral, definimos campos de vetores ao longo de curvas:
\begin{definition}
    \label{loc:2.9_:_def_4_:_campo_ao_longo_de_uma_curva}
    Se $c:I\to M$ é uma curva parametrizada, uma aplicação $V$, que leva cada $t \in I$ a um vetor $V(t) \in T_{c(t)}M$ é um campo de vetores ao longo da curva $c$.

    Dizemos que $V$ é diferenciável se, para todo $f \in \mathcal{D}$, a função $V(t)f$ é diferenciável em $I$.
\end{definition}
Perceba que, tomando $X$ um campo vetorial diferenciável em $M$, podemos definir $V(t)=X(c(t))$. Se isso acontecer, dizemos que $V$ é um campo induzido por $X$. A recíproca nem sempre é verdade, isto é, nem todo campo ao longo de uma curva é induzido. De fato, considere a curva diferenciável
\begin{equation*}
    c:(-\pi,\pi) \ni t \to (\cos t,\sin2t)\in \mathbb{R}^{2}
\end{equation*}
e o campo $c'(t)=(\sin t, 2\cos2t)$ ao longo de $c$.
Temos que $c\left( -\frac{\pi}{2} \right)=c\left( \frac{\pi}{2} \right)$, mas $c'\left(-\frac{\pi}{2} \right)\neq c'\left( \frac{\pi}{2} \right)$. Logo, não existe campo diferenciável $X \in \chi (\mathbb{R}^{2})$, tal que $c'=X\circ c$.

Note que, nesse caso, o campo ao longo da curva não é induzido pois assume diferentes valores no mesmo ponto (em tempos diferentes). E se forçarmos a curva a ser simples, isto é, $c$ injetora? Ainda é possível encontrar um contra-exemplo:
\begin{equation*}
    c:(-1, 1) \ni t \to (t^{2},t^{3}) \in \mathbb{R}^{2}
\end{equation*}
e $c'(t)=(2t, 3t^{2})$.
Aqui, a impossibilidade de encontrar um campo diferenciável que induz $c'$ se da pelo fato $c$ formar um bico em torno da origem, o que só acontece pois $c'(0)=0$. Com efeito, se $X \in \chi(\mathbb{R}^{2})$ é um campo diferenciável, tal que $c'(t)=X(c(t))$, escrevendo $X(x,y)=(P(x,y),Q(x,y))$, temos que $P(c(t))=2t$. Logo, derivando em ambos os lados,
\begin{equation*}
    \frac{\partial P}{\partial x}(c(t))\frac{\partial t^{2}}{\partial t}(t) + \frac{\partial P}{\partial y}(c(t))\frac{\partial t^{3}}{\partial t}(t)=2,
\end{equation*}
o que é um absurdo quando $t=0$.

\begin{figure}[htbp]
    \centering
    \caption{Gráficos das funções $c(t)=(t^{2},t^{3})$ e $c'(t)=(2t,3t^{2})$.}

    \begin{tikzpicture}

    % ============================================================
    % First curve: c(t) = (cos t, sin 2t)
    % ============================================================
    \begin{axis}[
            name=first,
            axis lines=center,
            axis equal,
            xlabel={$x$},
            ylabel={$y$},
            xmin=-1.3, xmax=1.3,
            ymin=-1.5, ymax=1.5,
            xtick={-1,0,1},
            ytick={-1,0,1},
            samples=300,
            width=7cm,
            height=7cm,
            grid=both,
            title={$c(t)=(\cos t,\sin 2t)$},
        ]

        % Curve
        \addplot[
            very thick,
            domain=-pi:pi,
        ] ({cos(deg(x))},{sin(2*deg(x))});

        % Derivative vectors
        \foreach \t in {
        -2.7,-2.2,-1.7,-1.2,-0.7,-0.2,
        0.3,0.8,1.3,1.8,2.3,2.8
        } {
        \pgfmathsetmacro{\cx}{cos(deg(\t))}
        \pgfmathsetmacro{\cy}{sin(2*deg(\t))}
        \pgfmathsetmacro{\dx}{-sin(deg(\t))}
        \pgfmathsetmacro{\dy}{2*cos(2*deg(\t))}

        % Normalize derivative
        \pgfmathsetmacro{\norm}{sqrt(\dx*\dx+\dy*\dy)}
        \pgfmathsetmacro{\tx}{\dx/\norm}
        \pgfmathsetmacro{\ty}{\dy/\norm}

        % Endpoint of tangent vector
        \pgfmathsetmacro{\vex}{\cx+0.18*\tx}
        \pgfmathsetmacro{\vey}{\cy+0.18*\ty}

        % EXPLANATION: \edef expands the variables immediately into their numerical values
        % before handing the \draw command over to pgfplots to queue.
        \edef\temp{\noexpand\draw[-{Stealth}, thick] (axis cs:\cx,\cy) -- (axis cs:\vex,\vey);}
        \temp
        }

    \end{axis}

    % ============================================================
    % Second curve: c(t) = (t^2, t^3)
    % ============================================================
    \begin{axis}[
            at={(first.east)},
            anchor=west,
            xshift=2cm,
            axis lines=center,
            axis equal,
            xlabel={$x$},
            ylabel={$y$},
            xmin=-0.3, xmax=1.3,
            ymin=-1.3, ymax=1.3,
            xtick={0,1},
            ytick={-1,0,1},
            samples=300,
            width=7cm,
            height=7cm,
            grid=both,
            title={$c(t)=(t^2,t^3)$},
        ]

        % Curve
        \addplot[
            very thick,
            domain=-1:1,
        ] ({x^2},{x^3});

        % Derivative vectors
        \foreach \t in {
        -0.9,-0.7,-0.5,-0.3,-0.1,
        0.1,0.3,0.5,0.7,0.9
        } {
        \pgfmathsetmacro{\cx}{\t*\t}
        \pgfmathsetmacro{\cy}{\t*\t*\t}

        \pgfmathsetmacro{\dx}{2*\t}
        \pgfmathsetmacro{\dy}{3*\t*\t}

        % Normalize derivative
        \pgfmathsetmacro{\norm}{sqrt(\dx*\dx+\dy*\dy)}
        \pgfmathsetmacro{\tx}{\dx/\norm}
        \pgfmathsetmacro{\ty}{\dy/\norm}

        % Endpoint of tangent vector
        \pgfmathsetmacro{\vex}{\cx+0.18*\tx}
        \pgfmathsetmacro{\vey}{\cy+0.18*\ty}

        % Use the \edef trick here as well
        \edef\temp{\noexpand\draw[-{Stealth}, thick] (axis cs:\cx,\cy) -- (axis cs:\vex,\vey);}
        \temp
        }

    \end{axis}

\end{tikzpicture}

    \vspace{6pt}
    \small Fonte: Elaborado pelo autor (2026).
    \label{fig:graficos_funcoes}
\end{figure}

Ajustamos então nossa pergunta: Se $c$ é uma curva simples e regular, isto é, injetora, com derivada não nula, e $V$ um campo ao longo de $c$, existe um campo $X$, tal que $V=X \circ c$? A resposta agora é positiva, se assumirmos que $c$ está definida em um intervalo compacto (incentivamos o leitor a encontrar um contraexemplo quando uma curva é simples e regular, porém definida em um intervalo não fechado). Precisamos primeiro extender a noção de vetor derivada para curvas em variedades diferenciáveis:
\begin{definition}
    \label{loc:2.9_:_def_5_:_campo_velocidade_da_curva}
    O campo velocidade de uma curva parametrizada é o campo
    \begin{equation*}
        \frac{dc}{dt}(t) :=dc_{t}\left( \frac{d}{dt}(t) \right),
    \end{equation*}
    onde $\frac{d}{dt}(t)=1$, na identificação do espaço tangente de $I$ com $\mathbb{R}$.
\end{definition}
Agora, conseguimos enunciar o resultado:
\begin{theorem}
    \label{loc:teorema_:_extensão_de_um_campo_ao_longo_de_curva_regular_simples_para_um_campo_na_variedade}
    Seja $c:[a,b]\to M$ uma curva parametrizada regular simples, isto é, com $\frac{dc}{dt}(t)\neq0$, para todo $t \in I$ e $c$  injetora. Se $V$ é um campo vetorial diferenciável ao longo de $c$, existe $Y \in \chi(M)$, tal que $V=Y \circ c$.
    \begin{proof}

    \end{proof}
\end{theorem}
Esse resultado será útil posteriormente, quando reduziremos noções envolvendo campos em uma variedade apenas para campos ao longo de curvas.

Com a noção de métrica e de curva, podemos definir comprimentos na variedade:
\begin{definition}
    \label{loc:2.9_:_def_6_:_comprimento_da_curva}
    Se $c:I\to M$ é uma curva parametrizada e $c|_{[a,b]}$ um segmento, definimos o comprimento do segmento como
    \begin{equation*}
        l_{a}^{b}(c)=\int_{a}^{b}\left<\frac{dc}{dt}(t),\frac{dc}{dt}(t)\right>_{c(t)}^{\frac{1}{2}}dt.
    \end{equation*}
\end{definition}
Em breve, mostraremos certas curvas interessantes que minimizam localmente o comprimento. Tais curvas, denominadas geodésicas, são de grande interesse no nosso estudo.

O próximo resultado, de cunho técnico, exibirá que toda variedade diferenciável Hausdorff, de base enumerável admite uma métrica riemanniana.
\begin{theorem}
    \label{loc:2.10_:_prop_2_:_existência_de_métrica}
    Seja $M$ uma variedade diferenciável Hausdorff e de base enumerável (\Cref{loc:axiomas_sobre_topologia_das_variedades}). Então, $M$ é uma variedade riemanniana, isto é, existe uma métrica riemmanama em $M$.
    \begin{proof}

        Pelo \Cref{loc:2.10_:_teo_3_:_condição_para_partição_da_unidade}, $M$ possui partição da unidade, isto é, existe uma coleção $\{ f_{\alpha} \}$, com $f_{\alpha}:M\to \mathbb{R}$, tal que, $f_{\alpha}\geq0$ e $suppf_{\alpha}\subset V_{\alpha}=\mathbf{x}_{\alpha}(U_{\alpha})$, para todo $\alpha$. Mais ainda, a coleção $\{ V_{\alpha} \}$ é localmente finita e $\sum_{\alpha}f_{\alpha}(p)=1$, para todo $p \in M$.

        Definimos $\left<\cdot,\cdot\right>^{\alpha}$ em cada $V_{\alpha}=\mathbf{x}_{\alpha}(U_{\alpha})$, pondo, para $p \in V_{\alpha}$,
        \begin{equation*}
            \left<u,v\right>_{p}^{\alpha}=\left<(d\mathbf{x}_{\alpha})_{p}^{-1}(u),(d\mathbf{x}_{\alpha})_{p}^{-1}(v)\right>,
        \end{equation*}
        em que o produto interno da direita é o produto interno euclidiano, identificando $T_{\mathbf{x}_{\alpha}^{-1}(p)}U_{\alpha}$ com $\mathbb{R}^{n}$, conforme feito em \Cref{loc:!resultados_extras_importantes!identificar_espaco_tangente}

        Mostremos que isso define uma métrica riemmaniana em cada $V_{\alpha}$:
        A simetria e o fato de ser positiva definida seguem do produto interno euclidiano. Para mostrar a diferenciabilidade, tomemos uma parametrização $\mathbf{y}$ em uma vizinhança de $p$. Escrevendo $\frac{\partial}{\partial y_{i}}(q)=d\mathbf{y}_{q}(0, \dots, 0, 1, 0, \dots, 0)$, temos que
        \begin{equation*}
            \left<\frac{\partial}{\partial y_{i}}(q),\frac{\partial}{\partial y_{j}}(q)\right>_{\mathbf{y}(q)}^{\alpha}=\left<d\mathbf{x}_{\alpha}^{-1} \circ d\mathbf{y}_{q}\left(\frac{\partial}{\partial y_{i}}(q)\right),d\mathbf{x}_{\alpha}^{-1} \circ d\mathbf{y}_{q}\left(\frac{\partial}{\partial y_{j}}(q)\right)\right>,
        \end{equation*}
        o qual varia diferencialmente em $q$, pelos axiomas de variedade diferenciável.

        Definimos agora uma métrica global em $M$, pondo
        \begin{equation*}
            \left<u,v\right>_{p}=\sum_{\alpha}f_{\alpha}(p)\left<u,v\right>_{p}^{_{\alpha}}.
        \end{equation*}
        É facil ver que tal aplicação é simétrica e positiva definida. A diferenciabilidade segue do fato da soma e produto de funções diferenciáveis ser diferenciável.
    \end{proof}
\end{theorem}
\section{Conexões Afins}
\label{loc:capítulo_3_:_métricas_e_conexões.conexões_afins}
Nessa seção daremos um passo pra trás e analizaremos outra ideia da topologia diferencial: conexões afins.
Já vimos campos vetoriais e curvas em variedades, junto com funções diferenciáveis e suas derivadas. Também podemos pensar em como uma função varia ao longo de um campo vetorial. Existe, entretanto, uma técnica do cálculo diferencial que ainda não está clara nas variedades: como podemos somar e subtrair vetores em diferentes espaços tangentes? Ou melhor, como derivar um campo vetorial em relação a outro campo vetorial?

Para fixar as ideias, considere uma variedade $M$, pontos $p,q \in M$ e seus espaços tangentes $T_{p}M$ e $T_{q}M$. Buscamos uma maneira de identificar vetores de $T_{p}M$ com vetores de $T_{q}M$.
Um palpite para fazer isso é, se $p, q$ pertencem a mesma vizinhança parametrizada, escrever um vetor $v \in T_{p}M$ como $v=\sum v_{i}X_{i}(p)$ e o identificar com o vetor $u=\sum v_{i}X_{i}(q) \in T_{q}M$. Essa identificação é um isomorfismo de espaços vetoriais, porém não é independente da parametrização. Outra forma, mais interessante, é levar em consideração a trajetória que o vetor percorre na variedade. Mais precisamente, dados $p,q \in M$, pensemos em uma transformação $\Gamma(\alpha, v) = u$, que recebe um caminho diferenciável $\alpha$ ligando $p$ a $q$ e um vetor $v \in T_{p}M$ e devolve um vetor $u \in T_{q}M$. Exigimos também que
\begin{enumerate}
    \item $\Gamma(\alpha \circ \beta, v) = \Gamma(\beta, \Gamma(\alpha, v));$
    \item $\Gamma(\alpha, \cdot):T_{p}M\to T_{q}M$ seja isometria de espaços vetoriais;
    \item $\Gamma(\alpha|_{[t_{0},t]},v)$ varie diferenciávelmente em $t$.
\end{enumerate}

$\Gamma$ é chamado de transporte paralelo de $v$ ao longo de $\alpha$.

Esta abordagem, de grande valor histórico, apesar de intuitiva, introduz algumas dificuldades quando realizamos cálculos. Por exemplo, não é simpes definir e demonstrar propriedades acerca de derivadas de campos vetoriais ao longo de outros campos.

Dito isso, seguiremos a tradição moderna e introduziremos a noção de conexão, a qual é equivalente a noção de transporte paralelo, de uma forma que será explicitada.
\begin{definition}
    \label{loc:2.1_:_def_1_:_conjunto_dos_campos_de_vetores}
    Seja $M$ uma variedade diferenciável. Definimos $\chi(M)$ como o conjunto de todos os campos de vetores em $M$ de classe $C^{\infty}$.
\end{definition}
\begin{definition}
    \label{loc:2.1_:_def_2_:_conexão_afim.definição}
    Uma conexão afim em uma variedade $M$ é uma aplicação
    \begin{equation*}
        \nabla:\chi(M)\times \chi(M) \ni (X, Y) \rightsquigarrow \nabla_{X}Y \in \chi(M),
    \end{equation*}
    que, sendo $f, g \in \mathcal{D}(M)$ e $X, Y, Z \in \chi(M)$, satisfaz as seguintes propriedades:
    \begin{enumerate}
        \item $\nabla_{fX+gY}Z=f\nabla_{X}Z + g\nabla_{Y}Z$;
        \item $\nabla_{X}(Y+Z)=\nabla _XY+\nabla_{X}Z$;
        \item $\nabla_{X}(fY)=f\nabla_{X}Y+(Xf)Y$.
    \end{enumerate}

    Isto é, $\nabla$ é linear na primeira entrada, distribui para a soma na segunda e satisfaz uma certa identidade para o produto.
\end{definition}
Intuitivamente, uma conexão é uma noção da derivada de um campo vetorial ao longo de outro campo vetorial. A propriedade $1$ é a linearidade da direção da derivada, a propriedade $2$ é análoga a regra da soma usual e a propriedade $3$ é análoga a regra do produto usual.  Apresentaremos adiante que toda variedade riemanniana admite uma única conexão afim que atende certas propriedades, o que remove a ambiguidade da escolha de uma conexão em uma variedade e realça a importância do estudo das conexões afins.
\begin{example}
    \label{loc:2.1_:_def_2_:_conexão_afim.exemplo}
    Quando $M=\mathbb{R}^{n}$, identificando $e_{i}$ com $X_{i}$ e, sendo $X = \sum\alpha_{i}e_{i}$ e $Y=\sum\beta_{j}e_{j}$ campos vetoriais diferenciáveis,
    \begin{equation*}
        \nabla_{X}Y=\sum_{j}\left( \sum_{i}\alpha_{i}\frac{\partial\beta_{j}}{\partial x_{i}} \right)e_{j}
    \end{equation*}
    é uma conexão afim. Mais ainda, sendo $p \in \mathbb{R}^{n}$ e $c:(-\epsilon,\epsilon):\mathbb{R}^{n}$ dada por $c(t)=X(p)t+p$, temos que
    \begin{equation*}
        (\nabla _{X}Y)_{p}=(Y \circ c)'(0)=\lim_{ t \to 0 } \frac{Y(c(t))-Y(p)}{t}.
    \end{equation*}
    De fato,
    \begin{align*}
        (Y \circ c)'(0) & = \left. \frac{d}{dt} Y(c_1(t), \dots, c_n(t)) \right|_{t=0}                            \\
                        & = \left. \sum_i \frac{\partial Y}{\partial x_i}(c(t)) \frac{dc_i}{dt}(t) \right|_{t=0}  \\
                        & = \sum_i \frac{\partial Y}{\partial x_i}(p) \frac{dc_i}{dt}(0)                          \\
                        & = \sum_i \left( \sum_j \frac{\partial \beta_j}{\partial x_i}(p) e_j \right) \alpha_i(p) \\
                        & = \sum_j \left( \sum_i \alpha_i(p) \frac{\partial \beta_j}{\partial x_i}(p) \right) e_j \\
                        & = (\nabla_{X}Y)_{p}.
    \end{align*}
    Em outras palavras, a conexão Levi-Civita de $\mathbb{R}^{n}$ coincide com a derivada usual.
\end{example}
Esse exemplo permite olhar para as conexões afins como uma generalização da ideia de derivada direcional no $\mathbb{R}^{n}$. Como é de se esperar, imitando a sua análoga para os espaços euclidianos, a conexão afim é uma propriedade local. Isto é, $(\nabla_{X}Y)_{p}$ depende apenas de $Y$ e do valor que $X$ assume em $p$:
\begin{lemma}
    \label{loc:lema_:_conexão_afim_é_uma_propriedade_local.}
    Seja $p \in M$ e $v \in T_{p}M$. Se $X, Y \in \chi(M)$ são tais que $X_{p}=Y_{p}$, então, para qualquer campo $Z \in \chi(M)$, temos que
    \begin{equation*}
        (\nabla_{X}Z)_{p}=(\nabla _{Y}Z)_{p}.
    \end{equation*}
    \begin{proof}

        Seja $\mathbf{x}$ uma parametrização de uma vizinhança de $p$. Temos que
        \begin{equation*}
            X=\sum x_i X_{i} \ \ \ \ \ \ \ \ \ \ \ Y=\sum y_{j}X_{j},
        \end{equation*}
        onde $X_{i}=\frac{\partial}{\partial x_{i}}$.
        Logo, utilizando as propriedades da conexão,
        \begin{align*}
            \nabla_{X}Y & = \nabla_{\sum x_{i}X_{i}}\left( \sum y_{j}X_{j} \right)                       \\
                        & =\sum_{i} x_{i}\left( \nabla_{X_{i}}\left( \sum_{j} y_{j}X_{j} \right) \right) \\
                        & =\sum_{i,j}x_{i}\nabla_{X_{i}}(y_{j}X_{j})                                     \\
                        & =\sum_{i,j}x_{i}(y_{j}\nabla_{X_{i}}X_{J}+(X_{i}y_{j})X_{j}                    \\
                        & = \sum_{i,j}x_{i}y_{j}\nabla_{X_{i}}X_{j}+\sum_{i,j}x_{i}(X_{i}y_{j})X_{j}.
        \end{align*}
        Como, para cada par $i, j$, o campo $\nabla_{X_{i}}X_{j}\in \chi(M)$, podemos escrever
        \begin{equation*}
            \nabla_{X_{i}}X_{j}=\sum \Gamma_{i,j}^{k}X_{k},
        \end{equation*}
        onde os $\Gamma_{i,j}^{k}$ são funções reais diferenciáveis definidas na vizinhança coordenada.
        Dessa forma,
        \begin{align*}
            \nabla_{X}Y & = \sum_{i,j} x_{i}y_{j}\left( \sum_{k}\Gamma_{i,j}^{k}X_{k}\right)+\sum_{i,j}x_{i}(X_{i}y_{j})X_{j}   \\
                        & =\sum_{k}\left( \sum_{i,j}x_{i}y_{j}\Gamma_{i,j}^{k} \right)X_{k} + \sum_{i,j}x_{i}(X_{i}y_{j})X_{j}.
        \end{align*}
        Observe que, para todo $j$, temos $X(y_{j})=\left( \sum x_{i}X_{i} \right)y_{j}=\sum x_{i}(X_{i}y_{j})$. Assim, trocando os índices do segundo somatório e utilizando essa igualdade, segue que
        \begin{equation*}
            \nabla_{X}Y=\sum_{k}\left( \sum_{i,j}x_{i}y_{j}\Gamma_{i,j}^{k} \right)X_{k} + \sum_{k}X(y_{k})X_{k}= \sum_{k}\left[\left( \sum_{i,j}x_{i}y_{j}\Gamma_{i,j}^{k} \right)+X(y_{k})\right]X_{k}.
        \end{equation*}
        Em outras palavras, $(\nabla_{X}Y)_{p}$ depende apenas de $x_{i}(p)$, $y_{k}(p)$ e das derivadas $X(y_{k})(p)$.

        Escrevendo então $Z=\sum z_{i}X_{i}$, temos que
        \begin{align*}
            (\nabla_{X}Z)_{p} & = \left.\sum_{k}\left[\left( \sum_{i,j}x_{i}z_{j}\Gamma_{i,j}^{k} \right)+X(z_{k})\right]X_{k}\right|_{p} \\
                              & =\left.\sum_{k}\left[\left( \sum_{i,j}y_{i}z_{j}\Gamma_{i,j}^{k} \right)+Y(z_{k})\right]X_{k}\right|_{p}  \\
                              & = (\nabla_{Y}Z)_{p}.
        \end{align*}
    \end{proof}
\end{lemma}
Antes de enunciarmos o próximo resultado, mostremos outro lema técnico:
\begin{lemma}
    \label{loc:lema_:_todo_vetor_pode_ser_estendido_a_um_campo.lema}
    Seja $p \in M$ e $v \in T_{p}M$. Existe um campo vetorial $X \in \chi (M)$, tal que $X_{p}=v$.
    \begin{proof}

        Seja $\mathbf{x}:U\to V$ uma parametrização em torno de $p$ e $X_{i}=\frac{\partial}{\partial x_{i}}$. Podemos escrever
        \begin{equation*}
            v=\sum v_{i}X_{i}(p).
        \end{equation*}
        Note que podemos pensar em $v_{i}$ como uma função constante de $V$ em $\mathbb{R}$.
        Defina o campo $X$ em $V$, pondo
        \begin{equation*}
            X_{q}=\sum v_{i}(q)X_{i}(q).
        \end{equation*}
        Note que $X$ é diferenciável, pois $X_{i}$ e $v_{i}$ são diferenciáveis. Também, não é difícil ver que $X_{p}=v$.
        Podemos então estender $X$ de forma que esteja definido em toda variedade $M$. Para isso, tomemos $\bar{U}\subset U$ compacto, conforme feito em \Cref{loc:def_11_:_funções_bump}, de forma que $p \in \bar{U}$, e $f:M\to \mathbb{R}$ uma função bump, com suporte compacto em $\bar{U}$. O campo $\bar{X}={f}X$ é o que procuramos.
    \end{proof}
\end{lemma}
Tomando $p \in M$ e $v \in T_{p}M$ e utilizando o \Cref{loc:lema_:_todo_vetor_pode_ser_estendido_a_um_campo.lema} e o \Cref{loc:lema_:_conexão_afim_é_uma_propriedade_local.}, dado $Y \in \chi(M)$ podemos pensar em $(\nabla_{v}Y)_{p}$ substituindo $v$ por um campo $X \in \chi(M)$, com $X_{p}=v$. De fato,
\begin{equation*}
    (\nabla_{X}Y)_{p} = \left.\sum_{k}\left[\left( \sum_{i,j}x_{i}y_{j}\Gamma_{i,j}^{k} \right)+X(y_{k})\right]X_{k}\right|_{p} = \sum_{k}\left[\left( \sum_{i,j}v_{i}y_{j}\Gamma_{i,j}^{k} \right)+v(y_{k})\right]X_{k}.
\end{equation*}
Dessa observação, conseguimos enunciar e demostrar o resultado seguinte:
\begin{proposition}
    \label{loc:2.2_:_prop_1_:_derivada_covariante}
    Seja $M$ uma variedade diferencial, com uma conexão afim $\nabla$. Existe uma única correspondência que associa um campo vetorial $V$ ao longo de uma curva parametrizada $c:I\to M$ a outro campo $\frac{DV}{dt}$, chamado de derivada covariante, ao longo de $c$, tal que, sendo $V$ e $W$ campos ao longo de $c$ e $f:I\to \mathbb{R}$ :
    \begin{enumerate}
        \item $\frac{D}{dt}(V+W)= \frac{D}{dt}V +\frac{D}{dt}W$;
        \item $\frac{D}{dt}(fV)=\frac{df}{dt}V+f\frac{DV}{dt}$;
        \item Se $V = Y(c(t))$, com $Y \in \chi(M)$, então $\frac{DV}{dt}(t)=\nabla_{\frac{dc}{dt}(t)}Y(c(t))$.
    \end{enumerate}
    \begin{proof}

        Seguindo o \Cref{loc:lema_:_definição_local_e_unicidade_levam_a_definição_global}, mostraremos primeiro a unicidade e depois definiremos $\frac{D}{dv}$ em vizinhanças coordenadas.
        Suponha que exista $\frac{D}{dt}$ satisfazendo 1, 2 e 3. Seja $\mathbf{x}:U\to M$ uma parametrização, tal que $\mathbf{x}(U)\cap c(I)\neq \emptyset$ e seja $\tilde{c}(t)=\mathbf{x}^{-1} \circ c(t)=(x_{1}(t), \dots, x_{n}(t))$ a expressão em coordenadas de $c$ em um aberto de $I$ contido em $\mathbf{x}(U)$.
        Sejam $X_{i}=\frac{\partial}{\partial x_{i}}$ e $V$ um campo vetorial ao longo de $c$.
        Escrevendo $V=\sum v_{j}X_{j}$, temos, pelas propriedades da derivada covariante, que
        \begin{equation*}
            \frac{DV}{dt}=\frac{D}{dt}\left( \sum_{j}v_{j}X_{j} \right)=\sum_{j}\frac{D}{dt}(v_{j}X_{j})=\sum_{j}\frac{dv_{j}}{dt}X_{j}+ \sum_{j}v_{j}\frac{DX_{j}}{dt}.
        \end{equation*}
        Note que
        \begin{equation*}
            \frac{DX_{j}}{dt}=\nabla_{\frac{dc}{dt}}X_{j}=\nabla_{\sum_{i}\frac{dx_{i}}{dt}X_{i}}X_{j}=\sum\frac{dx_{i}}{dt}\nabla_{X_{i}}X_{j}
        \end{equation*}
        implicando que
        \begin{equation*}
            \frac{DV}{dt}=\sum_{j}\frac{dv_{j}}{dt}X_{j}+\sum_{i,j}v_{j}\frac{dx_{i}}{dt}\nabla_{X_{i}}X_{j} \tag{*}.
        \end{equation*}
        Tal expressão de $\frac{DV}{dt}$ depende apenas da expressão em coordenadas de $c$, do campo $V$ e da conexão $\nabla$ e, portanto, é única.

        Agora, definimos $\frac{DV}{dt}$ em toda vizinhança parametrizada utilizando a equação $(*)$. O resultado segue do \Cref{loc:lema_:_definição_local_e_unicidade_levam_a_definição_global}.
    \end{proof}
\end{proposition}
Esse resultado extende a noção de conexão para campos definidos ao longo de curvas diferenciáveis. Isso permite, por exemplo, que calculemos acelerações de curvas computando a derivada covariante do campo velocidade de uma curva $c$.
A terceira propriedade da derivada covariante afirma que, calcular a covariante de um campo induzido, nada mais é do que calcular a conexão afim desse campo ao longo da curva $c(t)$. A derivada covariante também é análoga a derivada covariante da Geometria Diferencial, com a única diferença sendo que na Geometria Riemanniana, precisamos da noção de conexão, a qual é dispensável nas superfícies diferenciáveis devido a presença do espaço ambiente.

Dando sequência, retomamos a noção de transporte paralelo:
\begin{definition}
    \label{loc:2.5_:_def_3_:_paralelismo}
    Seja $M$ uma variedade diferenciável com uma conexão afim $\nabla$. Um campo vetorial $V$ ao longo de uma curva $c$ é dito paralelo quando $\frac{DV}{dt}=0$.
\end{definition}
\begin{proposition}
    \label{loc:2.6_:_prop_2_:_transporte_paralelo}
    Sejam $M$ uma variedade diferenciável com uma conexão afim $\nabla$, $c:I\to M$ uma curva parametrizada e $V_{0} \in T_{c(t_{0})}M$. Então, existe um único campo de vetores paralelo $V$ ao longo de $c$, tal que $V(t_{0})=V_{0}$.
    \begin{proof}

        Caso particular: $c(I)$ está contido em uma vizinhança parametrizada $\mathbf{x}(U)$.
        Sejam $\tilde{c}(t)=\mathbf{x}^{-1} \circ c(t) = (x_{1}(t), \dots, x_{n}(t))$ e $V_{0}=\sum V_{0}^{j}X_{j}$, onde $X_{j}=\frac{\partial}{\partial x_{j}}(c(t_{0}))$.
        Pela demonstração de \Cref{loc:2.2_:_prop_1_:_derivada_covariante}, temos que um campo $V$ em $\mathbf{x}(U)$ ao longo de $c$, com $V(t_{0})=V_{0}$ é paralelo se, e somente se,
        \begin{equation*}
            0=\frac{DV}{dt}=\sum_{j}\frac{dv^{j}}{dt}X_{j}+\sum_{i,j}\frac{dx_{i}}{dt}v^{j}\nabla_{X_{i}}X_{j}.
        \end{equation*}
        Fazendo $\nabla_{X_{i}}X_{j}=\sum_{k}\Gamma_{i,j}^{k}X_{k}$, temos que
        \begin{equation*}
            0=\sum_{j}\frac{dv^{j}}{dt}X_{j}+\sum_{i,j,k}\frac{dx_{i}}{dt}v^{j}\Gamma_{i,j}^{k}X_{k}.
        \end{equation*}
        Trocando o índice do primeiro somatório para $k$, segue que
        \begin{equation*}
            0=\sum_{k}\left\{  \frac{dv^{k}}{dt}+\sum_{i,j}v^{j}\frac{dx_{i}}{dt}\Gamma_{i,j}^{k}  \right\}X_{k},
        \end{equation*}
        isto é, para todo $k=1, \dots, n$,
        \begin{equation*}
            0=\frac{dv^{k}}{dt}+\sum_{i,j}v^{j}\frac{dx_{i}}{dt}\Gamma_{i,j}^{k}. \tag{*}
        \end{equation*}
        Dessa forma, $V$ é paralelo se, e somente se satisfaz a equação $(*)$, para todo $k=1, .., n$.

        Note que, para todo $k$,
        \begin{equation*}
            \frac{dv^{k}}{dt}=\begin{bmatrix}
                \sum_{i}\Gamma_{i,1}^{k}\frac{dx_{i}}{dt} & \cdots & \sum_{i}\Gamma_{i,n}^{k}\frac{dx_{i}}{dt}
            \end{bmatrix}\begin{bmatrix}
                v_{1} \\ \vdots \\ v_{n}
            \end{bmatrix}.
        \end{equation*}
        Logo,
        \begin{equation*}
            \begin{bmatrix}
                \frac{dv^{1}}{dt} \\ \vdots \\ \frac{dv^{n}}{dt}
            \end{bmatrix} = \begin{bmatrix}
                \sum_{i}\Gamma_{i,1}^{1}\frac{dx_{i}}{dt} & \cdots & \sum_{i}\Gamma_{i,n}^{1}\frac{dx_{i}}{dt} \\
                \vdots                                    & \ddots & \vdots                                    \\
                \sum_{i}\Gamma_{i,1}^{n}\frac{dx_{i}}{dt} & \cdots & \sum_{i}\Gamma_{i,n}^{n}\frac{dx_{i}}{dt}
            \end{bmatrix} \begin{bmatrix}
                v_{1}  \\
                \vdots \\
                v_{n},
            \end{bmatrix},
        \end{equation*}
        isto é, o sistema de equações é linear e, portanto, possuí uma única solução. Em outras palavras, existe um único campo $V$ que satisfaz as condições desejadas. \textbf{PLACEHOLDER: CITAR ALGUM LIVRO QUE CONTÉM ESSA DEMONSTRAÇÃO}

        Caso geral:
        Seja $c:I\to M$ curva parametrizada. Fixando $t_{0}\in I$, para todo $t_{1} \in I$ tal que $t_{0} \leq t_{1}$, o intervalo $[t_{0}, t_{1}]$ é compacto. Logo, como $c$ é contínua, $c([t_{0},t_{1}])$ é compacto isto é, podemos cobrir o segmento com finitas vizinhanças coordenadas, digamos $\mathbf{x}_{i}(U_{i})$ com $i=1, \dots, m$.
        Pelo que foi argumentado no caso particular, existe um único $V_{i}$ paralelo ao longo de $c$ em cada $\mathbf{x}_{i}(U_{i})$. O resultado segue aplicando o \Cref{loc:lema_:_definição_local_e_unicidade_levam_a_definição_global} no segmento $c([t_{0},t_{1}])$.
    \end{proof}
\end{proposition}
O transporte paralelo, conforme já comentado, identifica diferentes espaços tangentes através de um caminho na variedade. É possível mostrar que as noções de conexão afim e transporte paralelo são equivalentes observando que
\begin{equation*}
    (\nabla_XY)_{p}=\lim_{ t \to 0 } \frac{\Gamma(\alpha|_{[0,t]},Y(\alpha(t)))^{-1}-Y(p)}{t}
\end{equation*}
com $\alpha$ curva diferenciável tal que $\alpha(0)=p$ e $\alpha'(t)=X(\alpha(t))$, é, de fato, uma conexão afim e seu transporte paralelo é dado por $\Gamma$. O leitor interessado pode consultar \textbf{PLACEHOLDER: ACHAR ALGUM LIVRO QUE TEM ISSO AQUI DEMONSTRADO}.
Note a semelhança desta fórmula com a derivada de campos vetoriais usual do $\mathbb{R}^{n}$.

Dada uma variedade diferencial com partição de unidade, existem diversas conexões afins nesta. De fato, seja $\{ (U_{\alpha}, \mathbf{x}_{\alpha}) \}$ uma estrutura diferencial de uma variedade $M$, de forma que os $V_{\alpha}=\mathbf{x}_{\alpha}(U_{\alpha})$ formem uma família localmente finita. Em cada $V_{\alpha}$ podemos definir uma conexão $\nabla^{\alpha}$ tomando uma matriz $M^{\alpha}_{k}=[\Gamma_{i,j}^{\alpha,k}]$ e pondo $\nabla^{\alpha}_{X_{i}}X_{j}=\sum\Gamma_{i,j}^{\alpha,k}X_{k}$. Definimos uma conexão afim global pondo $(\nabla_{X}Y)_{p}=\sum_{\alpha}f_{\alpha}(p)(\nabla_{X}^{\alpha}Y)_{p}$, conforme feito na demonstração do \Cref{loc:2.10_:_prop_2_:_existência_de_métrica}. Devido a isso, o comportamento do transporte paralelo é imprevisível. Entretanto, o transporte paralelo assume uma forma muito mais intuitiva na conexão de Levi-Civita, a qual será introduzida adiante.

Agora, voltaremos a falar de variedades riemannas, buscando relacionar o conceito de métrica com o de conexão afim.
\begin{definition}
    \label{loc:3.1_:_def_4_:_conexão_compatível}
    Seja $M$ uma variedade Riemanniana de métrica $\left<,\right>$ e conexão $\nabla$.
    A conexão é dita compatível com a métrica quando, para toda curva parametrizada $c$ e quaisquer pares de campos de vetores paralelos $P$ e $P'$ ao longo de $c$ tivermos
    \begin{equation*}
        \left<P,P'\right> = \text{constante}.
    \end{equation*}
\end{definition}
Essa noção generaliza a regra de Leibnz para produtos internos na geometria diferencial. A próxima proposição torna isso evidente e motiva tal definição:
\begin{proposition}
    \label{loc:3.2_:_prop_3_:_leibniz_para_conexão_compatível}
    Uma conexão $\nabla$ em uma variedade Riemanniana com uma métrica $\left<,\right>$ é compatível se, e somente se, para todo $V,W$ campos diferenciáveis ao longo da curva parametrizada $c:I\to M$ tivermos
    \begin{equation*}
        \frac{d}{dt}\left<V,W\right> =\left<\frac{DV}{dt},W\right> +\left<V,\frac{DW}{dt}\right>. \tag{*}
    \end{equation*}
    \begin{proof}

        Se vale $(*)$, então tomando dois campos $P$ e $P'$ paralelos ao longo de uma curva parametrizada $c$, temos que
        \begin{equation*}
            \frac{d}{dt}\left<P,P'\right>=0+0 = 0.
        \end{equation*}
        Como tal função está definida em um intervalo, o qual é conexo, temos que $\left<P,P'\right>$ é constante.
        Reciprocamente, seja $\{ P_{1}(t_{0}), \dots, P_{n}(t_{0}) \}$ uma base ortonormal de $T_{c(t_{0})}M$ em relação ao produto interno $\left<,\right>_{c(t_{0})}$.
        Utilizando \Cref{loc:2.6_:_prop_2_:_transporte_paralelo}, extendemos cada $P_{i}(t_{0})$ a um campo $P_{i}(t)$ paralelo ao longo de $c$.
        Da hipótese, segue que $\left<P_{i}(t),P_{j}(t)\right>=\left<P_{i}(t_{0}),P_{j}(t_{0})\right>=\delta_{i,j}$, isto é, $\{ P_{1}(t), \dots, P_{n}(t) \}$ é uma base ortonormal de $T_{c(t)}M$, para todo $t \in I$.
        Podemos então escrever $V=\sum v_{i}P_{i}$ e $W=\sum w_{j}P_{j}$, com $v_{i}=\left<V,P_{i}\right>$ e $w_{j}=\left<W,P_{j}\right>$.
        Segue que
        \begin{equation*}
            \frac{DV}{dt}=\frac{D\left( \sum v_{i}P_{i} \right)}{dt}=\sum\frac{Dv_{i}P_{i}}{dt}=\sum \frac{dv_{i}}{dt}P_{i}+vi\frac{DP_{i}}{dt} = \sum\frac{dv_{i}}{dt}P_{i},
        \end{equation*}
        notando que $\frac{DP_{i}}{dt}=0$.
        De forma análoga, é possível mostrar que
        \begin{equation*}
            \frac{DW}{dt}=\sum\frac{dw_{j}}{dt}P_{j}.
        \end{equation*}
        Assim,
        \begin{align*}
            \left<\frac{DV}{dt},W\right> + \left<V,\frac{DW}{dt}\right> & = \sum_{i}\frac{dv_{i}}{dt}\left<P_{i},W\right> + \sum_{i}\frac{dw_{i}}{dt}\left<V,P_{i}\right> \\
                                                                        & = \sum_{i} \frac{dv_{i}}{dt}w_{i} + \frac{dw_{i}}{dt}v_{i}                                      \\
                                                                        & = \frac{d}{dt}\left( \sum_{i}v_{i}w_{i} \right)                                                 \\
                                                                        & = \frac{d}{dt}\left<V,W\right>  ,
        \end{align*}
        em que a última desigualdade vale devido a expressão de $V$ e $W$ na base ortonormal.
    \end{proof}
\end{proposition}
Outra equivalência da compatibilidade é dada pelo corolário seguinte:
\begin{corollary}
    \label{loc:3.3_:_corolário_1_:_compatível_se_e_só_se_vale_leibniz_direcional}
    Uma conexão $\nabla$ em uma variedade Riemanniana $M$ munida com uma métrica $\left<,\right>$ é compatível com a métrica se, e somente se, para quaisquer $X, Y, Z \in \chi(M)$,
    \begin{equation*}
        X\left<Y,Z\right> =\left<\nabla_{X}Y,Z\right> +\left<Y,\nabla_{X}Z\right> . \tag{$*$}
    \end{equation*}
    \begin{proof}

        Suponha que $\nabla$ é compatível com a métrica. Sejam $p \in M$ e $c:I\to M$ uma curva parametrizada, com $c(t_{0})=p$ e $\frac{dc}{dt}(t_{0})=X_{p}$. Temos que
        \begin{align*}
            X_{p}\left<Y,Z\right> & =\left.\frac{d}{dt}\left<Y(c(t)),Z(c(t))\right>\right|_{t=t_{0}}                                                                                                                                               \\
                                  & =\left<\left.\frac{D(Y \circ  c)}{dt}(t),Z(c(t))\right> +\left<Y(c(t)),\frac{D(Z \circ c)}{dt}(t)\right> \right|_{t=t_{0}}                   & \text{pela [[3.2 - Prop 3 - Leibniz para conexão compatível]]}. \\
                                  & =\left<\frac{D(Y \circ  c)}{dt}(t_{0}),Z_{p}\right> +\left<Y_{p},\frac{D(Z \circ c)}{dt}(t_{0})\right>                                                                                                         \\
                                  & = \left<\left( \nabla_{\frac{dc}{dt}(t_{0})}Y \right)_{p},Z_{p}\right> +\left<Y_{p},\left( \nabla_{\frac{dc}{dt}(t_{0})}Z \right)_{p}\right> & \text{pela definição de conexão}.
        \end{align*}
        Logo, como $\frac{dc}{dt}(t_{0})=X_{p}$, pelo \Cref{loc:lema_:_todo_vetor_pode_ser_estendido_a_um_campo.lema}, temos que
        \begin{equation*}
            X_{p}\left<Y,Z\right> =\left<(\nabla_{X}Y)_{p},Z_{p}\right> + \left<Y_{p},(\nabla _{X}Z)_{p}\right> =(\left<\nabla_{X}Y,Z\right> +\left<Y,\nabla_{X}Z\right> )_{p},
        \end{equation*}
        que é o que procuramos.

        Reciprocamente, sejam $c:I\to M$ uma curva parametrizada e $V, W$ campos diferenciáveis ao longo de $c$.

        Note que se $X,Y,Z$ são campos definidos em um aberto $U$ de $M$, então vale $(*)$ para $X,Y,Z$. De fato, tomando $p \in U$, por \Cref{loc:def_11_:_funções_bump}, existe uma função diferenciável $f:M\to [0,1]$ com suporte compacto contido em $U$ e contendo $p$. Definimos então $\overline{X} =\overline{Y} =\overline{Z}=0$, se $q \not\in U$ e $\overline{X}=fX, \overline{Y}=fY, \overline{Z}=fZ$ caso contrário. Aplicamos $(*)$ nestes campos o que torna evidente que a igualdade vale em $p$ para os campos originais. Como $p$ é arbitrário, temos o desejado.

        Note também que, utilizando a definição de vetor tangente, a hipótese, e a relação da covariante com campos induzidos temos que
        \begin{align*}
            \left.\frac{d}{dt}\left<X_{i} \circ  c,X_{j} \circ  c\right>\right|_{t=t_{0}} & = \frac{dc}{dt}(t_{0})\left<X_{i},X_{j}\right>                                                                                \\
                                                                                          & = \left<\nabla_{\frac{dc}{dt}(t_{0})}X_{i},X_{j}\right> + \left<X_{i},\nabla_{\frac{dc}{dt}(t_{0})}X_{j}\right>               \\
                                                                                          & = \left.\left<\frac{D(X_{i} \circ  c)}{dt},X_{j}\right> + \left<X_{i},\frac{D(X_{j} \circ  c)}{dt}\right>  \right|_{t=t_{0}}.
        \end{align*}

        Agora, escreva $V=\sum v_{i}(X_{i}\circ c)$ e $W=\sum w_{j}(X_{j} \circ c)$. Temos que
        \begin{align*}
            \left.\left<\frac{DV}{dt},W\right> +\left<V,\frac{DW}{dt}\right> \right|_{t=t_{0}} & = \sum_{i,j}v_{i}w_{j}\left<\frac{D(X_{i} \circ  c)}{dt},X_{j}\right> + \sum_{i,j}v_{i}w_{j}\left<X_{i},\frac{D(X_{j} \circ  c)}{dt}\right> \\
                                                                                               & =\sum_{i,j}v_{i}w_{j}\frac{d}{dt}\left<X_{i} \circ  c,X_{j} \circ  c\right>                                                                 \\
                                                                                               & =\frac{d}{dt}\left<\sum_{i}v_{i}(X_{i} \circ  c),\sum_{j}w_{j}(X_{j} \circ  c)\right>                                                       \\
                                                                                               & =\frac{d}{dt}\left<V,W\right>,
        \end{align*}
        que é o que queríamos mostrar.
    \end{proof}
\end{corollary}
Intuitivamente, esse corolário diz que a regra de Leibniz também vale quando derivamos o produto interno de dois campos vetoriais ao longo de outro campo vetorial.

É evidente que buscamos sempre uma conexão afim compatível com a métrica. Outra propriedade interessante de uma conexão afim é a simetria:
\begin{definition}
    \label{loc:3.4_:_def_5_:_conexão_simétrica..definição}
    Uma conexão afim $\nabla$ em uma variedade diferenciável $M$ é dita simétrica quando

    \begin{equation*}
        \nabla_{X}Y-\nabla_{Y}X=\left[X,Y\right]
    \end{equation*}
    para todo $X, Y \in \chi(M)$.
\end{definition}
Podemos caracterizar as conexões afins simétricas analisando seu comportamento nos vetores induzidos por uma parametrização.
\begin{proposition}
    \label{loc:3.4_:_def_5_:_conexão_simétrica..proposição_:_equivalência}
    Em uma parametrização $\mathbf{x}:U\to M$, uma conexão afim $\nabla$ é simétrica se, e só se,
    \begin{equation*}
        \nabla_{X_{i}}X_{j}-\nabla_{X_{j}}X_{i} = \left[X_{i}, X_{j}\right]=X_{i}(X_{j})-X_{j}(X_{i})=0,
    \end{equation*}
    para todo $i\neq j$.
    \begin{proof}

        Se $\nabla$ for simétrica, a implicação é trivial. Agora, supondo que vale a igualdade, sejam $X,Y \in \chi(M)$. Escrevendo $X=\sum\alpha_{i}X_{i}$ e $Y=\sum\beta_{j}X_{j}$, temos que
        \begin{equation*}
            \nabla _{X}Y=\nabla_{\left( \sum\alpha_{i}X_{i} \right)}\left( \sum \beta_{j}X_{j} \right)=\sum_{i,j}\alpha_{i}(\beta_{j}\nabla_{X_{i}}X_{j}+(X_{i}\beta_{j})X_{j}).
        \end{equation*}
        Logo,
        \begin{align*}
            \nabla_{X}Y-\nabla_{Y}X & =\sum_{i,j}\alpha _{i}(\beta_{j}\nabla_{X_{i}}X_{j}+(X_{i}\beta_{j})X_{j})-\beta_{j}(\alpha_{i}\nabla_{X_{j}}X_{i}+(X_{j}\alpha_{i})X_{i}) \\
                                    & = \sum_{i.j}\alpha_{i}\beta_{j}(\nabla_{X_{i}}X_{j}-\nabla_{X_{j}}X_{i})+\alpha_{i}(X_{i}\beta_{j})X_{j}-\beta_{j}(X_{j}\alpha_{i})X_{i}   \\
                                    & = \sum_{i,j}\alpha_{i}(X_{i}\beta_{j})X_{j}- \sum_{i,j}\beta_{j}(X_{j}\alpha_{i})X_{i},
        \end{align*}
        implicando que, para qualquer $f \in \mathcal{D}(M)$,
        \begin{align*}
            (\nabla_{X}Y-\nabla_{Y}X)f & =\left( \sum_{i,j}\alpha_{i}(X_{i}\beta_{j})X_{j}- \sum_{i,j}\beta_{j}(X_{j}\alpha_{i})X_{i} \right)f                                                                                          \\
                                       & = \sum_{i,j}\alpha_{i}(X_{i}\beta_{j})X_{j}f- \sum_{i,j}\beta_{j}(X_{j}\alpha_{i})X_{i}f                                                                                                       \\
                                       & =\sum_{i,j}\alpha_{i}\frac{\partial\beta_{j}}{\partial x_{i}}\frac{\partial f}{\partial x_{j}} - \sum_{i,j}\beta_{j}\frac{\partial\alpha_{i}}{\partial x_{j}}\frac{\partial f}{\partial x_{i}} \\
                                       & = (XY-YX)f                                                                                                                                                                                     \\
                                       & = [X,Y]f,
        \end{align*}
        onde a penultima desigualdade segue trocando os índices do segundo somatório.
    \end{proof}
\end{proposition}
O nome da propriedade é justificado pelo seguinte corolário:
\begin{corollary}
    \label{loc:3.4_:_def_5_:_conexão_simétrica..corolário_1}
    Em uma parametrização $\mathbf{x}:U\to M$, uma conexão afim $\nabla$ é simétrica se, e só se, escrevendo $\nabla_{X_{i}}X_{j}=\sum \Gamma_{i,j}^{k}X_{k}$, temos, $\Gamma_{ij}^{k}=\Gamma_{ji}^{k}$.
    \begin{proof}

        Imediato da \Cref{loc:3.4_:_def_5_:_conexão_simétrica..proposição_:_equivalência}.
    \end{proof}
\end{corollary}
A \Cref{loc:3.4_:_def_5_:_conexão_simétrica..proposição_:_equivalência} nos ajuda a compreender melhor o comportamento de uma conexão simétrica. Fica claro que uma conexão simétrica preserva uma propriedade interessante do cálculo no $\mathbb{R}^{n}$: derivadas parciais comutam. Na terminologia das conexões, isso se torna $\nabla_{X_{i}}X_{j}=\nabla_{X_{j}}X_{i}$. Também, impor simetria sobre uma conexão impica automaticamente que a conexão é compatível com uma certa estrutura intrínsica da variedade diferenciável: o colchetes. Isto é, a ``não comutatividade" da conexão é exatamente a ``não comutatividade" medida pelo colchetes.

Agora, apresentaremos um dos resultados centrais desse capítulo: o Teorema de Levi-Civita. Esse reultado garante que, dada qualquer variedade riemanniana, existe exatamente uma conexão afim que é compatível com a métrica e simétrica.
\begin{theorem}[Levi-Civita]
    \label{loc:3.6_:_teorema_levi_civita}
    Dada $M$ uma variedade riemanniana, existe uma única conexão afim $\nabla$ em $M$, tal que:
    \begin{enumerate}
        \item $\nabla$ é simétrica;
        \item $\nabla$ é compatível com a métrica riemanniana.
    \end{enumerate}
    \begin{proof}

        Mostremos primeiro a unicidade.
        Suponha que $\nabla$ exista. Pela compatibilidade da métrica, para quaisquer $X, Y, Z \in \chi(M)$, temos que:

        \begin{align*}
            X\left<Y,Z\right>  & = \left<\nabla_{X}Y,Z\right> +\left<Y,\nabla_{X}Z\right> \tag1    \\
            Y\left<Z,X\right>  & = \left<\nabla_{Y}Z,X\right> + \left<Z,\nabla_{Y}X\right>  \tag2  \\
            Z \left<X,Y\right> & = \left<\nabla_{Z}X,Y\right> + \left<X,\nabla_{Z}Y\right>.  \tag3
        \end{align*}

        Calculando $(1)+(2)-(3)$ e utilizando a compatibilidade e a simetria da conexão, temos

        \begin{align*}
            X\left<Y,Z\right> +Y\left<Z,X\right> -Z\left<X,Y\right> & = \left<\nabla_{X}Y,Z\right> +\left<Y,\nabla_{X}Z\right> + \left<\nabla_{Y}Z,X\right> + \left<Z,\nabla_{Y}X\right>       \\
                                                                    & \quad - \left<\nabla_{Z}X,Y\right> - \left<X,\nabla_{Z}Y\right>                                                          \\
                                                                    & = \left<\nabla_{X}Y,Z\right> +\left<Y,\nabla_{X}Z\right> + \left<\nabla_{Y}Z,X\right> + \left<Z,\nabla_{Y}X\right>       \\
                                                                    & \quad - \left<\nabla_{Z}X,Y\right> - \left<X,\nabla_{Z}Y\right>+ \left<\nabla_{X}Y,Z\right> - \left<\nabla_{X}Y,Z\right> \\
                                                                    & =\left<\nabla_{X}Z,Y\right>-\left<\nabla_{Z}X,Y\right> +\left<\nabla_{Y}Z,X\right> -\left<\nabla_{Z}Y,X\right>           \\
                                                                    & \quad \left<\nabla_{X}Y,Z\right> - \left<\nabla_{Y}X,Z\right> +2\left<Z,\nabla_{Y}X\right>                               \\
                                                                    & = \left<\nabla_{X}Z-\nabla_{Z}X,Y\right> + \left<\nabla_{Y}Z-\nabla_{Z}Y,X\right>+\left<\nabla_{X}Y-\nabla_{Y}X,Z\right> \\
                                                                    & \quad +2 \left<Z,\nabla_{Y}X\right>                                                                                      \\
                                                                    & = \left<[X,Z],Y\right> + \left<[Y,Z],X\right> +\left<[X,Y],Z\right> +2 \left<Z,\nabla_{Y}X\right>.
        \end{align*}

        Resolvendo para $\left<Z,\nabla_{Y}X\right>$, segue que

        \begin{align*}
            \left<Z,\nabla_{Y}X\right> & =\frac{1}{2} \{X\left<Y,Z\right> +Y\left<Z,X\right> -Z\left<X,Y\right>                 \\
                                       & \quad - \left<[X,Z],Y\right> - \left<[Y,Z],X\right> -\left<[X,Y],Z\right>  \}. \tag{4}
        \end{align*}

        Por fim, resta observar que
        \begin{equation*}
            \nabla_{Y}X = \sum \left<P_{i},\nabla_{Y}X\right> P_{i},
        \end{equation*}
        onde $\{ P_{i} \}$ é uma base ortonormal de $T_{p}M$. Por $(4)$, temos que a expressão de $\nabla_{X}Y$ depende apenas de $X, Y$ e da base ortonormal escolhida, implicando na unicidade.

        Agora, para mostrar a existência, seja $\mathbf{x}:U\to V$ um sistema de coordenadas de $V \subseteq M$. Definimos, para cada par $X_{i}=\frac{\partial}{\partial x_{i}}$ e $X_{j}=\frac{\partial}{\partial x_{j}}$,
        \begin{equation*}
            \nabla_{X_{i}}X_{j} = \sum \left<P_{i},\nabla_{X_{i}}X_{j}\right> P_{i},
        \end{equation*}
        onde $P_{i}$ é uma base ortonormal e $\left<P_{i},\nabla_{X_{i}}X_{j}\right>$ denota a expressão em $(4)$. O resultado segue do \Cref{loc:lema_:_definição_local_e_unicidade_levam_a_definição_global}.
    \end{proof}
\end{theorem}
De agora em diante, sempre que considerarmos uma variedade riamanniana, assumiremos que a variedade está munida com a conexão compatível com a métrica e simétrica. Tal conexão é chamada de conexão riemanniana ou conexão Levi-Civita.

Chamaremos os coeficientes $\Gamma_{i,j}^{k}$ de uma conexão Levi-Civita de símbolos de Christoffel.

Perceba a grandesa desse resultado: ele afirma que existe uma conexão que, de certa maneira, é intrínseca a variedade riemanniana. Quando consideramos a conexão Levi-Civita, não estamos fazendo uma conexão arbitrária entre planos tangentes, mas sim, estabelecendo um referencial que permite estudar propriedades de tal variedade riemannniana.
Podemos calcular os símbolos de Christoffel em uma parametrização apenas utilizando os produtos internos desta. De fato, escrevendo $g_{i,j}=\left<X_{i},X_{j}\right>$, temos que
\begin{equation*}
    \left<X_{l},\nabla_{X_{i}}X_{j}\right> = \left<X_{l},\sum_{k}\Gamma_{i,j}^{k}X_{k}\right>= \sum_{k}\Gamma_{i,j}^{k}g_{l,k}  .
\end{equation*}
Isso implica que
\begin{equation*}
    \begin{bmatrix}
        \left<X_{1},\nabla_{X_{i}}X_{j}\right> \\
        \vdots                                 \\
        \left<X_{n},\nabla_{X_{i}}X_{j}\right>
    \end{bmatrix} =
    M
    \begin{bmatrix}
        \Gamma_{i,j}^{1} \\
        \vdots           \\
        \Gamma_{i,j}^{n}
    \end{bmatrix},
    $$
        com
    $$
    M = \begin{bmatrix}
        g_{1,1} & \cdots & g_{1,n} \\
        \vdots  &        & \vdots  \\
        g_{n,1} & \cdots & g_{n,n}
    \end{bmatrix}.
    $$
        Como $M$ é uma matriz de Gram, ela é invertível. Digamos que
        $$
            M^{-1} = \begin{bmatrix}
                g^{1,1} & \cdots & g^{1,n} \\
                \vdots  &        & \vdots  \\
                g^{n,1} & \cdots & g^{n,n}
            \end{bmatrix}.
        $$
        Segue que
    $$
    \begin{bmatrix}
        \Gamma_{i,j}^{1} \\
        \vdots           \\
        \Gamma_{i,j}^{n}
    \end{bmatrix}=
    M^{-1}
    \begin{bmatrix}
        \left<X_{1},\nabla_{X_{i}}X_{j}\right> \\
        \vdots                                 \\
        \left<X_{n},\nabla_{X_{i}}X_{j}\right>
    \end{bmatrix}
\end{equation*}
e, portanto,
\begin{equation*}
    \Gamma_{i,j}^{k}=\sum_{l}\left<X_{l},\nabla_{X_{i}}X_{j}\right>g^{k,l}. \tag{*}
\end{equation*}
A partir da equação 4 do \Cref{loc:3.6_:_teorema_levi_civita},
\begin{align*}
    \left<X_l,\nabla_{X_i}X_j\right> & = \frac{1}{2} \Big\{ X_j\left<X_i,X_l\right> + X_i\left<X_l,X_j\right> - X_l\left<X_j,X_i\right>    \\
                                     & \quad - \left<[X_j,X_l],X_i\right> - \left<[X_i,X_l],X_j\right> - \left<[X_j,X_i],X_l\right> \Big\} \\
                                     & =\frac{1}{2}\Big\{X_{j}g_{i,l}+X_{i}g_{l,j}-X_{l}g_{j,i}\Big\}. \tag{**}
\end{align*}
De $(*)$ e $(**)$ segue que
\begin{equation*}
    \Gamma_{i,j}^{k}=\frac{1}{2}\sum_{l}\Big\{\frac{\partial}{\partial x_{j}}g_{i,l}+\frac{\partial}{\partial x_{i}}g_{l,j}-\frac{\partial}{\partial x_{l}}g_{j,i}\Big\}g^{k,l}.
\end{equation*}
Conforme comentado anteriormente, analisemos exemplos de conexões Levi-Civita e de transporte paralelo.
\begin{examples}
    \label{loc:exemplos_:_conexão_levi:civita_e_transporte_paralelo.exemplos}
    \begin{enumerate}
        \item A conexão definida no \Cref{loc:2.1_:_def_2_:_conexão_afim.exemplo} é Levi-Civita. Com efeito,
    \end{enumerate}
    \begin{equation*}
        \nabla_{e_{i}}e_{j}=\sum_{l}\left( \sum_{k}\delta_{k,i}\frac{\partial\delta_{l,j}}{\partial x_{k}} \right)e_{l} = \sum_{l}\frac{\partial\delta_{l,j}}{\partial x_{i}}e_{l}=\frac{\partial1}{\partial x_{i}}e_{j}=0,
    \end{equation*}
    isto é, todos os símbolos de Christoffel da conexão são nulos e, portanto, a conexão é simétrica. Para mostrar que é compatível com a métrica, observe que a derivada covariante, utilizando a expressão obtida na demonstração da \Cref{loc:2.2_:_prop_1_:_derivada_covariante}, coincide com a derivada usual e, portanto, vale a regra de Leibniz, isto é, a conexão é compatível com a métrica.
    2. Sendo $S \subseteq \mathbb{R}^{3}$ uma superfície regular e $X,Y$ campos de vetores em $S$, definimos
    \begin{equation*}
        (\nabla_{X}Y)_{p}=proj_{T_{p}M}(D_{X}Y(p)),
    \end{equation*}
    em que $proj$ é a projeção ortogonal e $D$ é a derivada usual de $\mathbb{R}^{3}$. É possível mostrar que essa é a conexão Levi-Civita da superfície e seus símbolos de Christoffel em uma paraetrização coincidem com os estudados na Geometria Diferencial.
    3. Pelo exemplo 2, conseguimos pensar no transporte paralelo estabelecido pela conexão Levi-Civita de maneira análoga a que fazemos na Geometria Diferencial. Em outras palavras, o transporte paralelo de um vetor ao longo da curva $c$ é, do ponto de vista de um habitande da variedade, uma forma de movimentar o vetor de um ponto até o outro de maneira que eles permanecam paralelos conforme ilustrado na imagem. \textbf{PLACEHOLDER: IMAGEM DO TRANSPORTE PARALELO NA ESFERA, FAZENDO ANALOGIA DO CARRO}.
\end{examples}
Esse capítulo deixa claro a importância do estudo das métricas riemannianas e das conexões e como esses conceitos carregam propriedades semelhantes as que temos no cálculo no $\mathbb{R}^{n}$, incluíndo propriedades de derivação em produtos internos e derivadas de campos vetoriais ao longo de outros campos.